\chapter{Studiu bibliografic}\label{ch:studiubib}

\pagestyle{fancy}

Studiul bibliografic a fost organizat pe tematici relevante pentru dezvoltarea unui instrument local de adnotare colaborativă. Referințele sunt grupate progresiv: de la instrumente și seturi de date existente, spre arhitecturi de segmentare și modele vizual-lingvistice, continuând cu tehnici de adnotare asistată și pattern-uri de interfață, și încheind cu mecanisme de sincronizare, calitate, trasabilitate și confidențialitate.

% ---------------------------------------------------------------
\section{Instrumente de adnotare semantică}

În raportul tehnic din 2017 al lui Breitenfeld și Muller-Birn se realizează o comparație a instrumentelor de adnotare semantică utilizate în contexte de cercetare, cu accent pe legarea conținutului la vocabulare controlate și baze de cunoștințe~\cite{Breitenfeld2017}. Autorii analizează uneltele după tipul de utilizator, formatele suportate, sursele de cunoștințe (RDF/SKOS), modul de interacțiune și contextul de rulare (web, desktop sau infrastructuri de cercetare). Concluzia principală este că majoritatea instrumentelor sunt orientate spre utilizatori tehnici, iar suportul pentru colaborare și trasabilitate rămâne limitat, ceea ce îngreunează adopția în echipe interdisciplinare. Deși lucrarea vizează adnotarea semantică și nu pe cea vizuală, observațiile privind uzabilitatea, colaborarea și trasabilitatea sunt direct relevante pentru proiectul de față, în care aceleași cerințe apar în procesul de etichetare a imaginilor.

% ---------------------------------------------------------------
\section{Seturi de date și benchmark-uri}

PASCAL Visual Object Classes (VOC) reprezintă unul dintre primele benchmark-uri utilizate pe scară largă pentru detecție și segmentare~\cite{10.1007/s11263-009-0275-4}. Setul de date include etichete pentru 20 de clase comune, bounding box-uri pentru detecție și măști pentru segmentare. Competițiile anuale organizate în jurul PASCAL VOC au consolidat metrici standard, precum Average Precision (AP) pentru detecție și IoU/mean IoU pentru segmentare, și au furnizat un cadru consistent de comparare a metodelor, de la HOG/SVM până la arhitecturi bazate pe rețele convoluționale profunde. Rămâne relevant ca reper metodologic și punct de referință în analiza evoluției adnotărilor.

Microsoft COCO extinde semnificativ PASCAL VOC, propunând un set de date de referință pentru detecție, segmentare și descriere de imagini în scene complexe cu obiecte multiple~\cite{lin2015microsoftcococommonobjects}. Colecția include sute de mii de imagini și peste două milioane de instanțe adnotate, organizate pe zeci de clase și supercategorii. Anotările sunt detaliate (segmentare pe instanță, bounding box-uri și descrieri textuale), iar evaluarea se bazează pe metrici de precizie calculate la multiple praguri IoU, pentru a reflecta mai realist calitatea localizării. COCO a standardizat protocoalele de evaluare și a devenit referința principală pentru compararea modelelor în sarcini de adnotare și segmentare.

Torralba și Efros analizează bias-ul seturilor de date din recunoașterea vizuală și arată că modelele învață adesea caracteristici specifice colecțiilor de antrenare, nu concepte vizuale generale~\cite{5995347}. Autorii compară mai multe seturi de date (PASCAL VOC, ImageNet, LabelMe, SUN09, Caltech101, MSRC) prin experimente de tip „Name That Dataset" și teste de generalizare încrucișată. Rezultatele indică o scădere semnificativă a acurateței atunci când antrenarea și testarea se realizează pe surse diferite, iar bias-ul de selecție, de captură și cel al exemplelor negative afectează direct robustețea modelelor. Deși studiul precede era rețelelor profunde și se bazează pe metode HOG/BoW cu SVM, concluziile privind bias-ul rămân actuale și susțin importanța controlului distribuției datelor și a fluxurilor de asigurare a calității.

% ---------------------------------------------------------------
\section{Arhitecturi de segmentare}

U-Net propune o arhitectură complet convoluțională pentru segmentare semantică în condiții cu puține date adnotate~\cite{ronneberger2015unetconvolutionalnetworksbiomedical}. Structura în formă de U combină o cale de contractare pentru extragerea caracteristicilor cu una de expansiune pentru reconstrucția rezoluției, legate prin conexiuni de tip skip care păstrează detaliile spațiale. Modelul utilizează o strategie overlap-tile pentru procesarea imaginilor mari, augmentare cu deformări elastice și o funcție de pierdere ponderată care accentuează separarea obiectelor adiacente. Evaluările pe datele ISBI au confirmat eficiența abordării chiar și cu seturi de antrenare reduse. Pentru proiect, U-Net susține ideea de auto-segmentare rapidă în interfață, antrenare eficientă cu puține exemple și procesare pe dale pentru imagini de rezoluție mare.

Hariharan et al.\ propun un cadru unificat pentru detecție și segmentare la nivel de instanță, care prezice simultan bounding box-uri și măști pentru fiecare obiect~\cite{hariharan2014simultaneousdetectionsegmentation}. Metoda extinde detectoarele bazate pe regiuni prin adăugarea predicției de mască per instanță, oferind o înțelegere mai fină a scenei față de segmentarea semantică clasică. Relevanța pentru proiect este directă: adnotarea de instanțe presupune atât localizarea obiectelor, cât și delimitarea precisă a contururilor, ceea ce justifică fluxuri care combină desenarea de bounding box-uri cu rafinarea ulterioară a măștilor.

Kirillov et al.\ unifică segmentarea semantică și cea de instanțe într-o sarcină unică numită segmentare panoptică, în care fiecare pixel primește o categorie semantică și, unde este cazul, un identificator de instanță~\cite{kirillov2019panopticsegmentation}. Autorii introduc metrica Panoptic Quality (PQ), care combină calitatea segmentării cu cea a recunoașterii, oferind un criteriu de evaluare unificat. Lucrarea a influențat arhitecturi ulterioare (Panoptic FPN, UPSNet, modele bazate pe transformere) și este relevantă pentru proiect prin cerința de adnotări complete și consistente la nivel de imagine.

Mask2Former propune o arhitectură universală pentru segmentare, bazată pe clasificarea măștilor și pe utilizarea de object queries care poate aborda simultan segmentarea semantică, de instanțe și panoptică~\cite{cheng2022maskedattentionmasktransformeruniversal}. Inovația principală este mecanismul de masked attention: fiecare query poate accesa doar regiunea prezisă anterior, reducând zgomotul de fundal și accelerând convergența. Decodorul folosește o strategie multi-scale cu runde succesive pe rezoluții diferite, iar pierderile sunt calculate pe un eșantion de puncte informative, limitând consumul de memorie la antrenare. Rezultatele arată performanță SOTA simultană pe COCO panoptic, COCO instance și ADE20K semantic, cu convergență în aproximativ 50 de epoci față de sute de epoci la arhitecturi anterioare. Pentru proiect, Mask2Former este relevant prin posibilitatea de fine-tuning eficient pe seturi create colaborativ și prin generarea de propuneri cu focalizare locală, utilă pentru uneltele de pre-selectare și corecție rapidă.

% ---------------------------------------------------------------
\section{Modele vizual-lingvistice}

CLIP introduce o metodă de învățare a reprezentărilor vizual-lingvistice prin antrenarea simultană a unui encoder vizual și a unuia textual pe perechi imagine-text la scară mare~\cite{radford2021learningtransferablevisualmodels}. Obiectivul contrastiv maximizează similaritatea pentru perechile corecte și o minimizează pentru cele incorecte, ceea ce permite clasificare zero-shot: etichetele țintă sunt formulate ca prompturi textuale (de tip „A photo of a \{label\}."), iar imaginea este clasată prin potrivire cu descrierea cea mai probabilă. CLIP se dovedește competitiv cu modele supervizate pe ImageNet fără antrenare specifică și mai robust la schimbări de distribuție (schițe, randări). Encoderul vizual CLIP a devenit un fundament pentru modelele de segmentare cu vocabular deschis. Pentru proiect, CLIP este relevant ca bază pentru interogări textuale flexibile în pre-adnotare și pentru sugestii rapide fără antrenare suplimentară, dar necesită verificare umană și control al bias-urilor în fluxul colaborativ.

Liang et al.\ abordează o limitare specifică a CLIP la segmentare: performanța modelului scade semnificativ pe imagini mascate, din cauza diferenței de domeniu și a tokenilor de fundal neinformativi~\cite{liang2023openvocabularysemanticsegmentationmaskadapted}. Soluția propusă adaptează CLIP pe regiuni mascate prin construirea unui set de date zgomotos din COCO Captions și prin tehnica Mask Prompt Tuning (MPT), care înlocuiește tokenii de fundal cu tokeni de prompt antrenabili. Modelul OVSeg obține rezultate SOTA în regim zero-shot pe mai multe benchmark-uri de segmentare cu vocabular deschis. Limitările includ influența claselor văzute asupra propunerilor de măști și decalajul dintre clasele văzute și cele nevăzute. Pentru proiect, tehnica susține interogări textuale libere, pre-adnotare automată din metadate și adaptare eficientă cu resurse reduse.

YOLOE propune un model unificat, open-set și în timp real pentru detecție și segmentare, integrând prompturi textuale, vizuale și un mod fără prompt în aceeași arhitectură~\cite{wang2025yoloerealtimeseeing}. Blocul RepRTA rafinează embedding-urile textuale la antrenare, apoi le re-parametrizează în capul de clasificare fără cost suplimentar la inferență. Blocul SAVPE separă caracteristicile semantice de cele specifice promptului vizual, agregându-le eficient. Modul prompt-free LRPC localizează toate obiectele fără un vocabular explicit, asociind ancorele relevante cu un vocabular extins. Rezultatele arată o reducere de circa 3 ori a costului de antrenare față de YOLO-World, câștiguri pe LVIS în regim zero-shot și viteză mai mare la inferență. Limitările includ dependența de un vocabular static în modul prompt-free și necesitatea generării de pseudo-măști cu SAM pentru antrenarea segmentării. Pentru proiect, YOLOE este relevant prin suportul multimodal, posibilitatea unui mod de auto-etichetare eficient și eliminarea costurilor suplimentare la rulare prin re-parametrizare.

LLaVA propune o arhitectură care combină un encoder vizual (CLIP) cu un LLM (Vicuna) printr-un strat de proiecție, permițând dialogul vizual-textual detaliat~\cite{liu2023visualinstructiontuning}. Pentru antrenament, autorii generează dialoguri om-AI condiționate de imagini cu ajutorul GPT-4. Modelul atinge aproximativ 85.1\% din performanța GPT-4 pe sarcini vizuale complexe și obține rezultate competitive pe ScienceQA. Limitările includ tendința la halucinații și dificultăți la citirea textului mic sau a elementelor subtile de interfață, din cauza rezoluției limitate a encoderului vizual. Pentru proiect, LLaVA este relevant prin capacitatea de a răspunde la întrebări despre capturi de ecran din interfața de adnotare, oferind asistență imediată utilizatorului fără a efectua acțiuni automate.

DistilBERT reduce dimensiunea modelului BERT cu 40\% prin distilarea cunoștințelor în faza de pre-antrenare, unde un model student mai mic imită distribuțiile de ieșire ale modelului profesor prin o funcție de pierdere compusă~\cite{sanh2020distilbertdistilledversionbert}. Rezultatul este un model cu inferență mai rapidă și consum redus de resurse, care păstrează o mare parte din capacitatea de înțelegere a limbajului. Pentru proiect, DistilBERT este relevant pentru rularea locală a sarcinilor de procesare a limbajului natural, precum clasificarea intențiilor sau sugerarea etichetelor, reducând latența și costurile față de apelul unui model mai mare.

% ---------------------------------------------------------------
\section{Tehnici de adnotare asistată}

Settles descrie în raportul său tehnic principiile și strategiile de active learning, arătând că un model poate obține performanțe ridicate cu mai puține etichete dacă selectează instanțele cele mai informative pentru adnotare~\cite{settles.tr09}. Sunt acoperite scenarii de eșantionare din flux și dintr-un rezervor, precum și strategii de interogare bazate pe incertitudine, comitete de modele sau reducerea erorii așteptate. În practică, active learning reduce efortul de etichetare, dar introduce provocări legate de costuri variabile, adnotatori zgomotoși și bias indus de selecție. Pentru un instrument colaborativ, aceste idei permit prioritizarea imaginilor sau a regiunilor cu cel mai mare aport informațional, evitând munca redundantă. Abordările de tip batch-mode active learning pot distribui sarcinile eficient între mai mulți adnotatori, iar mecanismele de re-etichetare pot compensa erorile umane.

Tehnicile Grad-CAM folosesc gradientul din ultimul strat convoluțional pentru a produce o hartă a regiunilor discriminative, care poate fi suprapusă peste imagine pentru a vizualiza zonele relevante pentru decizia modelului~\cite{Selvaraju_2019}. Pentru proiect, aceste tehnici sunt utile ca mecanisme de transparență și asigurare a calității: pot arăta de ce modelul a sugerat o etichetă sau o mască, pot evidenția erori de focalizare (modelul urmărește fundalul în loc de obiect) și pot sprijini detectarea bias-ului din date. Afișarea hărților ca overlay în interfața colaborativă facilitează validarea rapidă a sugestiilor și corectarea interactivă, reducând timpul de revizie.

Polygon-RNN generează adnotări de instanțe ca poligoane produse secvențial de un RNN peste caracteristici extrase de un CNN, iar utilizatorul poate corecta vârfurile în timpul generării~\cite{castrejon2017annotatingobjectinstancespolygonrnn}. Modelul prezice vârfuri pe o grilă discretizată și se oprește la un token special de final, transformând desenarea completă a unui contur într-un proces de ajustare a câtorva vârfuri. Autorii raportează accelerări semnificative față de adnotarea manuală, cu un acord ridicat exprimat prin IoU. Pentru proiect, această abordare susține un flux de asistență AI cu feedback uman, în care interfața trebuie să permită corecturi rapide și să integreze imediat rezultatul în forma finală.

Deep Extreme Cut (DEXTR) reduce efortul de adnotare prin generarea unei măști precise din doar patru puncte extreme (sus, jos, stânga, dreapta), codate ca hărți Gaussiene și oferite ca intrare suplimentară modelului~\cite{maninis2018deepextremecutextreme}. Utilizatorul indică rapid limitele obiectului, fără a trasa poligoane complete, iar modelul produce o mască inițială de calitate. Limitarea apare la obiecte puternic ocluzionate sau cu forme neregulate, unde punctele extreme pot fi ambigue și necesită corecții. Pentru proiect, DEXTR susține un flux de tipul „indicare, apoi corecție", în care masca inițială poate fi ajustată cu pensula sau prin repoziționarea punctelor extreme pentru re-inferență rapidă. Acest tip de workflow scade timpul de revizie și reduce bias-ul de oboseală.

ScribbleBox propune un flux în două etape pentru adnotare video: componenta Curve-VOT urmărește obiectul printr-o curbă temporală cu câteva puncte de control, iar Scribble-VOS rafinează măștile pe baza scribble-urilor propagate cu un GCN~\cite{chen2020scribbleboxinteractiveannotationframework}. Pe setul DAVIS2017, metoda atinge 88.92\% J\&F cu circa 9.1 click-uri per secvență și corecții pe doar 4 cadre. Autorii raportează latențe detaliate: propagarea măștilor la 512px ajunge la circa 295\,ms/cadru, la 256px circa 152\,ms, iar cu cache circa 82\,ms; propagarea scribble-urilor este de 15--28\,ms, iar curve fitting și interacțiunea box sunt sub 15\,ms. Limitarea principală este dependența de calitatea urmăririi inițiale — obiectele mici sau cu mișcări imprevizibile necesită mai multe puncte de control. Pentru un instrument colaborativ, structura în două etape permite împărțirea sarcinilor pe niveluri de expertiză și impune o interfață asincronă cu indicatori de progres, cache persistent și mod rapid/precis pentru a gestiona latențele ridicate.

% ---------------------------------------------------------------
\section{Evoluția pattern-urilor UI/UX în adnotare}

Evoluția interfeței de adnotare reflectă trecerea de la interacțiunea manuală la cea asistată. LabelMe permite editarea directă a poligoanelor (zoom, pan, adăugare/ștergere de noduri), oferind precizie ridicată, dar cu un cost semnificativ pentru măști complexe~\cite{5483185}. Instrumente moderne precum ScribbleBox și CVAT reduc interacțiunea la indicii brute (casete sau scribble-uri), iar modelul generează rapid măști de calitate cu feedback aproape instant~\cite{chen2020scribbleboxinteractiveannotationframework,cvat2025}. Modele eficiente precum YOLOE susțin acest feedback în regim open-vocabulary fără costuri mari la inferență~\cite{wang2025yoloerealtimeseeing}.

CVAT structurează munca de adnotare în proiecte, sarcini și job-uri, reducând conflictele și permițând paralelizarea activității~\cite{cvat2025}. Un flux în două etape similar cu ScribbleBox — casete brute urmate de rafinare — permite alocarea sarcinilor pe niveluri de expertiză: adnotatorii juniori generează rapid casete brute, iar cei seniori rafinează măștile complexe. Pattern-urile de tip issues/comentarii și corecțiile non-distructive prin scribble-uri reduc refacerile și accelerează verificarea. Practici suplimentare, precum consensul și honeypot-urile, pot fi adăugate pentru a măsura consistența fără a întrerupe fluxul de lucru.

% ---------------------------------------------------------------
\section{Sincronizare colaborativă}

Sun et al.\ analizează diferențele dintre OT (Operational Transformation) și CRDT (Commutative Replicated Data Type) printr-un cadru de transformare unificat~\cite{Sun2020OTCRDT}. Autorii separă starea vizibilă de utilizator de structurile interne de sincronizare și arată că și CRDT implică transformări indirecte cu costuri suplimentare, ceea ce poate afecta performanța în sisteme interactive. Concluziile sugerează că, pentru aplicații cu server central și colaborare în timp real, un model operațional bazat pe propagarea evenimentelor este mai predictibil și mai eficient decât un CRDT complex. Pentru un instrument de adnotare colaborativă, această perspectivă susține trimiterea operațiilor (creare/modificare/ștergere) în locul transmiterii stării complete.

Arhitecturile event-driven facilitează colaborarea în timp real prin transmiterea asincronă a evenimentelor peste conexiuni persistente, iar STOMP (Simple Text Oriented Messaging Protocol) oferă un protocol simplu, bazat pe text, pentru comunicare de tip publicare/abonare peste TCP sau WebSocket~\cite{stomp_spec_1_2}. Protocolul definește frame-uri standard (CONNECT, SEND, SUBSCRIBE, MESSAGE, ACK/NACK, DISCONNECT) și separă canalul de transport de logica de rutare. Pentru proiect, acest model susține broadcast instant al schimbărilor pe canvas, blocarea și editarea obiectelor și interoperabilitatea dintre frontend și backend.

% ---------------------------------------------------------------
\section{Calitatea datelor și trasabilitate}

Lucrarea „Counting on Consensus" oferă un ghid pentru alegerea metricilor de acord între adnotatori (IAA), subliniind că procentul brut de acord este înșelător în clasele dezechilibrate și recomandă metrici corectate pentru hazard precum Kappa sau Alpha~\cite{james2026countingconsensusselectingright}. Pentru etichete continue sau ordonate sunt relevante Kappa ponderată sau ICC, iar pentru adnotări structurate sunt utile măsuri de suprapunere (F1/Dice) și metrici de frontieră. Deși focusul lucrării este pe NLP, principiile se transferă direct la adnotarea de imagini: sistemul poate calcula IAA pe etichete și pe măști, poate semnala cazurile ambigue și poate impune raportarea intervalelor de încredere, sprijinind un pipeline de QA bazat pe consistență, fără a confunda fiabilitatea cu validitatea.

Standardul W3C PROV-DM oferă un model formal pentru descrierea modului în care sunt generate și modificate entitățile digitale~\cite{Moreau2013PROVDM}. Modelul introduce concepte de bază — entitate, activitate, agent — și relații precum \textit{wasGeneratedBy} și \textit{wasDerivedFrom}, care permit reconstruirea istoricului unei adnotări. Aceste mecanisme susțin auditarea modificărilor, reproducibilitatea și controlul calității, mai ales atunci când adnotările sunt produse atât de utilizatori, cât și de modele automate.

Metadatele structurate (proiecte, sarcini, etichete, istoricul modificărilor) sunt potrivite pentru un RDBMS precum PostgreSQL, care oferă tranzacții ACID, control concurent și indexare eficientă~\cite{postgresql_docs}. Fișierele media sunt stocate în sisteme de fișiere sau object storage și sunt referențiate prin chei în baza de date, asigurând o separare clară între date structurate și cele nestructurate. În proiect, PostgreSQL este utilizat în principal pentru autentificare, dar rămâne opțiunea optimă pentru extensii viitoare care vor necesita organizarea metadatelor la scară mai mare.

% ---------------------------------------------------------------
\section{Scalabilitate și confidențialitate}

Neurosurgeon propune o schemă de inteligență colaborativă cloud-edge, în care o rețea neuronală este împărțită la un punct optim între dispozitivul utilizatorului și serverul cloud~\cite{10.1145/3037697.3037698}. Modelul estimează timpul de execuție și consumul de energie pe fiecare strat, atât local, cât și în cloud, și alege stratul la care activarea este transmisă către server, optimizând latența sau consumul energetic în funcție de lățimea de bandă disponibilă și de capacitățile dispozitivului. Evaluările pe mai multe arhitecturi arată că această împărțire reduce semnificativ latența sau energia față de rularea completă pe cloud sau exclusiv local. Pentru proiect, această abordare permite rularea locală a straturilor inițiale pentru feedback rapid (propuneri brute), cu delegarea straturilor costisitoare către backend pentru rafinare, menținând o experiență interactivă cu latență mică.

Lucrarea „Privacy-Preserving Collaborative Learning" propune o schemă de învățare colaborativă cu protecție eterogenă, în care participanții definesc bugete individuale de confidențialitate prin perturbarea cu zgomot Gaussian în procesul de agregare a modelelor~\cite{10163938}. Cadrul zCDP ofuscă atât modelele locale, cât și pe cel global, asigurând convergența către o soluție optimă și protejând datele față de server și față de ceilalți participanți. Conceptele din lucrare fundamentează politicile de autentificare și autorizare necesare într-un sistem de adnotare multi-tenant, asigurând izolarea datelor între adnotatori. Arhitectura susține un flux de pre-etichetare pe dispozitiv (on-device), în care serverul primește doar actualizări de model ofuscate, fără acces direct la imaginile brute ale utilizatorilor.